\chapter{Introduction and Context}%
\label{ch:intro}

Music notation is one of the most sophisticated symbolic languages ever devised: a
two-dimensional graphical system capable of encoding pitch, rhythm, dynamics, phrasing,
and expressive intent with sufficient precision to allow performers separated by centuries
to reproduce a composer's intentions from a printed page.
Yet the overwhelming majority of notated music exists today only as paper or
low-resolution scanned image, inaccessible to the computational tools that now drive
music education, music information retrieval (MIR), computational musicology, and
digital archiving.
\term{Optical Music Recognition} (OMR) is the discipline that bridges this divide: it
translates images of musical scores into structured, machine-readable representations in
the same way that Optical Character Recognition (OCR) transformed printed text into
searchable, editable data~\cite{shatriOPTICALMUSICRECOGNITION}.

This thesis describes the design, implementation, and quantitative evaluation of an
end-to-end OMR system tailored to a concrete and practically significant target:
the monophonic melody lines and chord symbols of jazz lead sheets, with a primary focus
on \textit{The Real Book}, among the most widely circulated fake-books in jazz
practice and pedagogy~\cite{kernfeldStoryFakeBooks2006}.

% ─────────────────────────────────────────────────────────────────────────────
\section{Motivation}
\label{sec:motivation}

\subsection{The Digitization Gap in Music}

Music archives, educational institutions, and individual musicians hold enormous
quantities of repertoire that exists only in paper or low-resolution scanned form.
Public libraries, conservatory collections, and private fake-book editions contain
hundreds of thousands of pages of notation that cannot be searched, transposed, played
back, or analysed without manual transcription into notation software.
Manual transcription is slow and demands specialist training that is unavailable at
scale---per-page effort varies widely by source complexity, but transcription labour is
repeatedly identified as a principal cost driver in corpus-level digitisation
projects~\cite{calvoZaragozaUnderstandingOMR2020, hajicOMRMetricsEvaluation2019}.

The practical consequences are significant.
A music educator who wishes to assign a jazz standard in a different key must either
own commercial software with the piece already encoded, transcribe it manually, or
distribute a photocopy that students cannot edit.
A musicologist studying harmonic patterns across a large corpus cannot apply statistical
methods to scanned PDFs without first converting them to machine-readable format.
A musician practising from a worn photocopy of \textit{The Real Book} cannot feed the
page to a practice app that transposes the melody to a B\(\flat\) instrument.

OCR technology began solving an equivalent problem for text decades ago, making printed books,
newspapers, and documents searchable and editable at almost zero marginal cost.
OMR is the musical analogue, but it remains far less mature: the visual language of
music notation is substantially more complex than alphabetic text, the available
annotated data is orders of magnitude smaller, and the community of researchers and
tools is correspondingly narrower---it is a more specialised field.

\subsection{Jazz Lead Sheets and \textit{The Real Book}}


Notated music exists in many formats---full orchestral scores, piano reductions,
chord charts, tablatures---in the same way written language ranges from novels to
tweets to legal documents.  Among these, jazz lead sheets occupy a particular
position.
A lead sheet is a compact score format that encodes only the most essential musical
information: a single-voice melody line written on a staff, chord symbols printed above
it indicating the harmonic context, and---occasionally---lyrics beneath the staff.
This format is the lingua franca of jazz practice: it allows musicians to memorise and
improvise over a tune from a minimal, portable representation without the burden of a
full orchestral score~\cite{reedHonorsLeadSheets2018}.

\textit{The Real Book} is the most widely circulated collection of jazz lead
sheets~\cite{martinezSevillaJazzLeadSheets2025}.
Originally assembled as an unauthorised compilation by Berklee College of Music students
in the 1970s, it circulated as a hand-copied samizdat for years before receiving
official publication.
It contains several hundred of the most important jazz standards---tunes by Miles Davis,
Thelonious Monk, Bill Evans, John Coltrane, and dozens of other canonical
figures---and is found on the music stands of jazz musicians at virtually every level of
practice, worldwide.

Despite its cultural centrality, the content of \textit{The Real Book} remains almost
entirely locked in paper or scanned-image form.
The officially published editions exist as PDFs, but those PDFs are image-based: the
musical content is not encoded as symbolic data, and standard music notation software
cannot interpret them directly.
Digitising \textit{The Real Book}---and the broader class of jazz lead sheets it
represents---would unlock a range of practical applications: automatic transposition for
B\(\flat\), E\(\flat\), and F instruments; instant MIDI playback for practice and
accompaniment; corpus-level harmonic analysis for musicologists; improved search and
cataloguing in digital music libraries; and accessible preparation of educational
materials for jazz educators.

\subsection{Why Existing OMR Tools Fall Short}

General-purpose OMR systems have advanced significantly in recent
years~\cite{shatriOPTICALMUSICRECOGNITION}, but they are designed and evaluated
primarily on classical full-score notation in standard engraving styles (Finale,
Sibelius, Dorico, or LilyPond's default font).
Lead sheets diverge from this baseline in three ways that, together, define the
specific challenges this thesis addresses.

\textbf{Engraving style.}
Many editions of \textit{The Real Book} use fonts and layouts that deliberately imitate
hand-lettered notation.  Models trained on classical-style corpora such as
PrIMuS~\cite{PrIMuSDataset} or DeepScoresV2~\cite{tuggenerDeepScoresV22020} do not
transfer naturally to this visual register.

\textbf{Chord symbols.}
Jazz chord symbols appear above the staff as free-form text and are a core part of
lead-sheet semantics (see~\Cref{sec:chord-symbols-primer} for the full grammar).
Standard OMR pipelines transcribe pitch and rhythm from the staff; they either ignore
chord symbols or treat them as noise.  Any lead-sheet OMR system must include a
complementary chord recogniser coordinated with melody transcription.

\textbf{Absence of labelled data.}
To the best of the author's knowledge---and consistent with recent surveys of OMR
datasets~\cite{shatriOPTICALMUSICRECOGNITION}, which list no such resource---no publicly
available, labelled dataset of Real Book pages with ground-truth symbolic transcriptions
exists as of 2026.
Supervised training of modern deep models without such a corpus requires synthetic
data generation and domain adaptation, whose calibration directly governs the
performance of the resulting system.  These gaps are surveyed in detail in
\Cref{ch:soa}.

This thesis assumes basic familiarity with Western staff notation---the staff and
treble clef, pitch and accidentals, note durations, time signatures, chord symbols, and
the lead-sheet format.
Readers without a musical background will find a self-contained primer on these elements
in \Cref{app:notation-primer}; the vocabulary design, data pipeline, and error analysis
of the following chapters assume that material.

% ─────────────────────────────────────────────────────────────────────────────
\section{Problem Formulation}
\label{sec:problem-formulation}

The goal of this thesis is to build a system that, given a single-page PDF or raster
image of a jazz lead sheet, automatically produces a structured, machine-readable transcription of its
contents: the sequence of notes, rests, and rhythmic markers making up the melody, and
the chord symbols specifying the harmonic progression.
Despite decades of OMR research, reliable automated transcription of real-world scores
remains an open problem~\cite{shatriOPTICALMUSICRECOGNITION}.
The challenge has several compounding sources of difficulty.

\textbf{Notation complexity.}
Unlike alphabetic text, where symbols are largely independent and arranged in a single
linear sequence, musical notation encodes information jointly across two spatial
dimensions.
A note's pitch is determined by its vertical position on the staff \textit{relative to
the clef}, not by the symbol shape alone; its duration is determined by a combination
of head shape, stem presence, flag count, and dot presence.
Beams create joint dependencies across groups of notes; ties extend a single pitch
across a barline.
Any recognition system must capture these two-dimensional, context-dependent symbol
interactions without relying on explicit per-symbol segmentation at inference time.

\textbf{Domain gap.}
The most widely used annotated OMR datasets---including
PrIMuS~\cite{PrIMuSDataset} and DeepScoresV2~\cite{tuggenerDeepScoresV22020}---consist
of cleanly computer-engraved images rendered at high resolution with standard engraving
fonts.
Real-world documents, by contrast, are photographed under imperfect lighting
conditions, and suffer from scanner noise, perspective distortion, ink bleed, paper
yellowing, and JPEG compression artefacts.
A model trained exclusively on clean synthetic data generalises poorly to real scans:
the feature distribution shifts substantially when crisp black symbols on a white
background are replaced by blurred, uneven, noisy images with variable contrast.
Bridging this \term{domain gap} without real annotated training data requires a
scan-simulation augmentation pipeline that teaches the model to be invariant to
scan-related distortions.

\textbf{Domain specificity and data scarcity.}
As argued above, jazz lead sheets use an engraving style for which no publicly
labelled corpus currently exists.  Training a model for this domain therefore requires
constructing a synthetic dataset from scratch using a rendering pipeline tuned to the
jazz style, since no off-the-shelf labelled corpus is available.

\textbf{Dual-stream recognition.}
The lead-sheet format requires simultaneously processing two distinct visual streams
from the same image region: the melodic staff content, which calls for pattern-based
sequence recognition, and the chord symbols above the staff, which call for text OCR.
These two streams require different feature representations and must ultimately be
integrated into a coherent structured output.
A complete lead-sheet OMR system must solve both sub-problems and align their results.

This thesis addresses all four difficulties jointly: by constructing a synthetic
training corpus in the jazz engraving style, applying a calibrated scan-simulation
augmentation pipeline, training a CRNN-CTC model for melody recognition, and building
a separate OCR pathway for chord symbols that is fused into a unified output by the
inference pipeline.

% ─────────────────────────────────────────────────────────────────────────────
\section{Objectives}
\label{sec:objectives}

The \textbf{general objective} of this project is to develop, evaluate, and document a
proof-of-concept end-to-end OMR system for monophonic jazz lead sheets, and to
demonstrate that a CRNN-CTC architecture trained on synthetic data can bridge the
domain gap between clean rendered images and degraded scans---achieving an aggregate
Symbol Error Rate (SER) of at most \SI{5}{\percent} on a held-out scan-simulated test
split, so that the outcome can be assessed quantitatively in \Cref{ch:conclusions}.

The \textbf{specific objectives} are:

\begin{enumerate}
    \item Survey the state of the art in OMR, covering traditional modular pipelines,
          CRNN-CTC sequence models, and transformer-based end-to-end approaches, and
          justify the architectural choices made in this project.

    \item Design and implement a synthetic training dataset by rendering excerpts from
          the PrIMuS corpus~\cite{PrIMuSDataset} in a jazz engraving style using
          LilyPond and the LilyJAZZ font, paired with ground-truth annotations in the
          LMX token format~\cite{OMRResearchLmx2025}.

    \item Design and implement a scan-simulation augmentation pipeline calibrated
          against real Real Book scan characteristics, to close the domain gap at
          training time without requiring manually labelled real-world data.

    \item Design, implement, and train a CRNN-CTC model for monophonic staff-line
          recognition, and evaluate it rigorously using the SER~\cite{SymbolErrorRate}
          metric, targeting an aggregate SER of at most \SI{2}{\percent} on the clean
          test split and at most \SI{5}{\percent} on the scan-simulated test split.

    \item Perform quantitative error analysis to characterise systematic failure modes
          and guide iterative data and model refinements.

    \item Build a complete inference pipeline---encompassing PDF input handling, staff
          detection, melody recognition, grammar correction, and chord symbol
          extraction---that returns structured JSON output suitable for downstream
          musical applications.

    \item Quantify the synthetic-to-scan domain gap as the aggregate-SER difference
          between the clean and scan-simulated test splits, and reduce the residual gap
          on genuine scans by fine-tuning on a small set ($\approx$\,50 staves) of
          manually annotated Real Book pages, targeting at least a \SI{30}{\percent}
          relative SER reduction on those pages.
\end{enumerate}

% ─────────────────────────────────────────────────────────────────────────────
\section{Proposed Approach}
\label{sec:approach}

The core design choice of this thesis is a \term{CRNN-CTC} architecture: a
Convolutional Recurrent Neural Network trained with the Connectionist Temporal
Classification (CTC) loss.
This model class is well-matched to the left-to-right, variable-length token sequences
that arise from reading a single staff line from left to right: the convolutional
backbone extracts local visual features along the width of the staff image; two
bidirectional LSTM layers model long-range sequential dependencies; and CTC decoding
maps the resulting frame-level distributions to a token sequence without requiring
explicit segmentation of the image into individual symbols.

An alternative---transformer-based end-to-end models such as the Sheet Music
Transformer~\cite{riosVilaSheetMusicTransformer2024} or the OLIMPIC system for
pianoform music~\cite{mayerPracticalEndToEnd2024}---offers higher modelling capacity
but demands substantially larger labelled datasets and compute budgets than are
available within the scope of a Bachelor's thesis.
Traditional heuristic pipelines~\cite{dalitzComparativeStudyStaff2008} avoid data
requirements but accumulate segmentation errors that propagate irreversibly through the
recognition chain.
The CRNN-CTC architecture represents the best trade-off for this project's constraints;
the technical case is developed in \Cref{ch:design} and the
project-management criteria behind the decision are recorded in
\Cref{sec:mgmt-justification}.

The system is organised around a five-stage inference pipeline---preprocessing,
staff detection, melody recognition, grammar correction, and chord OCR---whose
outputs are merged into a structured JSON document.  The decomposition, its
contracts between stages, and its design rationale are presented in detail in
\Cref{ch:design}.

The training strategy also deserves mention.
In the absence of labelled real-world data, the model is first trained entirely on
synthetic LilyJAZZ-rendered images augmented with scan simulation, then fine-tuned on a
small annotated set of real Real Book pages.
This two-phase approach is a standard domain adaptation technique and is the practical
solution to the data scarcity problem described in \Cref{sec:problem-formulation}.

% ─────────────────────────────────────────────────────────────────────────────
\section{Scope}
\label{sec:scope}

The project is bounded by the following scope definition.

\textbf{In scope:}
\begin{itemize}
    \item Input handling for a single-page PDF or a single raster image; the first page
          of a multi-page PDF is processed.
    \item Monophonic (single-voice) melody lines in treble clef, as found in jazz lead
          sheets.
    \item Synthetic training data construction via LilyPond rendering of PrIMuS
          excerpts in the LilyJAZZ engraving style.
    \item Scan-simulation augmentation modelling noise, blur, perspective distortion,
          ink bleed, and photometric variation.
    \item Staff region detection as a preprocessing module.
    \item Design, training, and evaluation of a CRNN-CTC model for single-staff
          melody recognition.
    \item Chord symbol extraction via an OCR-based module.
    \item Quantitative evaluation using SER and qualitative error analysis.
    \item Domain-gap analysis and fine-tuning on a small set of annotated Real Book
          pages.
    \item A FastAPI web service exposing the full pipeline as a REST endpoint.
\end{itemize}

\textbf{Out of scope:}
\begin{itemize}
    \item Polyphonic transcription or simultaneous multi-voice recognition.
    \item Full-page layout analysis and reading-order recovery across multiple stave
          systems beyond the per-staff processing implemented.
    \item Handwritten scores (as opposed to printed or typeset notation).
    \item Non-Western notation systems or extended contemporary notation.
    \item Real-time or mobile-device deployment.
    \item Training large-scale transformer models from scratch.
    \item Multi-page PDF processing (only the first page of any multi-page document is
          read).
    \item Tuplets, repeat barlines, voltas, and navigation marks (D.C., D.S., segno,
          coda), which do not appear frequently enough in the in-scope repertoire to
          justify the vocabulary and grammar extensions they would require.
\end{itemize}

% ─────────────────────────────────────────────────────────────────────────────
\section{Document Structure}
\label{sec:structure}

The remainder of this document is organised as follows.

\Cref{ch:soa} surveys the state of the art in OMR, covering the progression
from heuristic modular pipelines through CRNN-CTC models to modern transformer-based
end-to-end architectures, and reviews the available datasets and symbolic
representations relevant to this work.

\Cref{ch:mgmt} covers project management: the justification of the chosen
architecture through comparative analysis, the work-breakdown structure and timeline,
the budget and economic analysis, and the project's sustainability and ethical impact
assessment.

\Cref{ch:design} describes the conceptual design of the system: the
pipeline decomposition, the data strategy, the LMX output representation and
vocabulary, the CRNN--CTC architectural choices, the two-stream approach for
melody and chord recognition, the training strategy, and the evaluation
design.

\Cref{ch:impl} describes the implementation: the technology stack, the
repository structure, the data and augmentation pipeline, the model and
training loop, the chord OCR network, the inference pipeline, the CLI and
web service, and the mechanisms used to keep the experiments reproducible.

\Cref{ch:results} presents experimental results: training curves, SER on the
synthetic test set, ablation studies, qualitative error analysis by failure category,
and domain-gap measurements with and without fine-tuning on real Real Book pages.

\Cref{ch:sustain} revisits sustainability and ethical considerations in light of
the completed work.

Finally, \Cref{ch:conclusions} draws conclusions, reflects on the extent to
which the stated objectives were achieved, and outlines directions for future work.
